\section{6. Simulating the sensor values (the physical model)}

BuildSim provides rooms but not their physics, so three things are simulated:
the people, the air in each room, and the devices that measure it.

\subsection{6.1 The people}

occupancysim builds each person's day from a seed and the date. Staff arrive
between 08:00 and 10:00, work in their own office, take fika around 09:45 and
14:30 with a probability of 70~\%, lunch between 12:00 and 13:00, and leave
between 16:00 and 19:00. Students attend one to three lectures in the slots
starting 08:15, 10:15, 13:15 and 15:15. Weekends are nearly empty, and three
guards patrol at night. A room's role follows from its area: an office holds
one person, a lecture room one per 2~m², a fika room one per 4~m².

One property of this model decides what a forecast can achieve. For each
lecture slot, occupancysim draws the floor's lecture rooms in random order and
books only as many as the students need. A given lecture room is therefore
usually empty and occasionally full, at times that do not repeat from day to
day. The fika rooms and the offices have a daily pattern; the lecture rooms do
not.

We changed occupancysim in two places. It now advances in fixed 5-second steps
where it used to step by whatever the wall clock gave it, and it records how
many people, and how many guards among them, are in every room at every whole
model minute (\code{GET /api/minutes}, IF-1). Before this, two runs of one day
saw the same plan through differently phased windows and disagreed in 0.3~\%
of samples. The day plans and everything occupancysim writes to BuildSim are
unchanged, and its own tests pass.

\subsection{6.2 The air in a room}

Each room is one well-mixed volume $V$ with $n$ people, an airflow $q$ and a
heating setpoint. The airflow and the setpoint are the states the room's damper
and radiator valve have \emph{reached}, which physics reads from BuildSim every
tick. This is where the loop closes: a command changes the reached state, the
reached state changes the next minute of physics, and the next reading shows
it.

\textbf{\COtwo{}} is a mass balance, with $C$ in ppm:
\[
V\,\frac{dC}{dt} = n\,G\cdot 10^{6} + q\,(C_{\mathrm{out}} - C),
\qquad
C_{\infty} = C_{\mathrm{out}} + \frac{n\,G\cdot 10^{6}}{q},
\qquad
\tau = \frac{V}{q}.
\]
People add \COtwo{} at $G$ per person; ventilation replaces room air with
outdoor air. Left alone, the room approaches $C_{\infty}$ with time constant
$\tau$.

\textbf{Temperature} is a heat balance over the room's thermal mass
$C_{\mathrm{th}}$:
\[
C_{\mathrm{th}}\,\frac{dT}{dt} = UA\,(T_{\mathrm{out}} - T)
 + \rho c_p\,q\,(T_{\mathrm{sup}} - T) + n\,P_{\mathrm{body}} + P_{\mathrm{rad}},
\qquad
P_{\mathrm{rad}} = P_{\max}\cdot\mathrm{clamp}\!\left(\frac{T_{\mathrm{set}} - T}{0.5\ \mathrm{K}},\,0,\,1\right).
\]
The terms are loss through the envelope, supply air arriving at
$T_{\mathrm{sup}}$, body heat, and the radiator, whose valve opens in
proportion to how far the room is below its setpoint.

Both equations are linear with the inputs held constant, so each step uses the
exact solution, $x(t+\Delta t) = x_{\infty} + (x(t) - x_{\infty})\,e^{-\Delta
t/\tau}$, and is stable for any step. Physics advances one model minute at a
time, with the people who were in the room at the start of the minute; a longer
gap, such as a clock jump, is cut into steps of at most one minute (FR-1).
Because the steps follow occupancysim's per-minute record and not physics' own
timing, the true state is the same function of model time in every run of a
day.

\begin{longtable}[]{@{}L{0.30\linewidth}L{0.22\linewidth}L{0.40\linewidth}@{}}
\toprule
\textbf{Constant} & \textbf{Value} & \textbf{Source or assumption} \\
\midrule
\endhead
Ceiling height & 2.7 m & Assumed; $V$ = floor area × height \\
Outdoor \COtwo{}, $C_{\mathrm{out}}$ & 420 ppm & \\
\COtwo{} per seated adult, $G$ & $5\cdot10^{-6}$ m³/s & About 18 l/h \\
Outdoor temperature, $T_{\mathrm{out}}$ & 8 °C & Luleå in autumn; constant \\
Supply air, $T_{\mathrm{sup}}$ & 18 °C & The air-handling unit heats outdoor air to this \\
Envelope loss, $UA$ & 1 W/K per m² floor & Assumed \\
Thermal mass, $C_{\mathrm{th}}$ & 100 kJ/K per m² floor & Air, furniture and inner walls; assumed \\
Body heat, $P_{\mathrm{body}}$ & 100 W per person & Seated adult \\
Radiator maximum, $P_{\max}$ & 60 W per m² floor & Assumed \\
Air, $\rho c_p$ & 1200 J/(m³·K) & \\
\bottomrule
\end{longtable}

\textbf{Energy proxy} (NFR-3). Physics reports three powers per room and
minute, and the evaluation sums them: fan power at 1.5~kW per m³/s of airflow
(a typical specific fan power), the heat needed to warm outdoor air to the
supply temperature, $\rho c_p\,q\,(T_{\mathrm{sup}} - T_{\mathrm{out}})$, which
is 12~kW per m³/s, and the radiator power from the heat balance.

\subsection{6.3 Ventilation sizing}

Each room's flows follow from its area and its design capacity, by the Swedish
workplace rule (AFS: 7~l/s per person plus 0.35~l/s per m²):
\[
q_{\mathrm{design}} = 7\ \mathrm{l/s}\times\text{capacity} + 0.35\ \mathrm{l/s\,m^2}\times\text{area}.
\]

\begin{longtable}[]{@{}L{0.22\linewidth}L{0.24\linewidth}L{0.46\linewidth}@{}}
\toprule
\textbf{Flow level} & \textbf{Value} & \textbf{Used for} \\
\midrule
\endhead
Maximum & 1.5 × $q_{\mathrm{design}}$ & The \COtwo{} override; the top of the damper's range \\
Design & $q_{\mathrm{design}}$ & The \code{constant} policy; the damper's fallback; degraded mode without \COtwo{} \\
Occupied minimum & 0.35 l/s per m² & Flow floor on weekdays 07--18 \\
Unoccupied minimum & 0.10 l/s per m² & Flow floor otherwise (building code, nobody present) \\
\bottomrule
\end{longtable}

The 50 rooms, as sized in the registry:

\begin{longtable}[]{@{}L{0.10\linewidth}L{0.07\linewidth}L{0.12\linewidth}L{0.11\linewidth}L{0.14\linewidth}L{0.12\linewidth}L{0.16\linewidth}@{}}
\toprule
\textbf{Role} & \textbf{Rooms} & \textbf{Area, m²} & \textbf{Capacity} & \textbf{Design flow, l/s} & \textbf{$\tau$ at design} & \textbf{\COtwo{} when full, at design} \\
\midrule
\endhead
Lecture & 20 & 63--237 & 31--119 & 239--916 & 12 min & about 1070 ppm \\
Fika & 3 & 100--214 & 25--53 & 210--446 & 21--22 min & about 1010 ppm \\
Office & 24 & 6--27 & 1 & 9--17 & 31--74 min & about 790 ppm \\
Unused & 3 & 8--26 & 1 & 10--16 & 38--73 min & \\
\bottomrule
\end{longtable}

Two things follow from the table. A full room at design flow settles right
around the 1000~ppm limit, which is why the maximum sits above design. And at
the occupied minimum every room has a time constant of about two hours
($2.7\ \mathrm{m} / 0.35\ \mathrm{l/s\,m^2}$), so a room that fills while its
damper sits at the floor is far behind. That is the lag a forecast is meant to
beat.

\subsection{6.4 The sensors and their faults}

Each sensor device samples its own room's truth from physics once per model
minute and reports what a sensor of its kind would:

\begin{longtable}[]{@{}L{0.16\linewidth}L{0.40\linewidth}L{0.36\linewidth}@{}}
\toprule
\textbf{Sensor} & \textbf{Noise} & \textbf{Reported as} \\
\midrule
\endhead
\COtwo{} & Gaussian with a standard deviation of half of (30~ppm + 3~\% of the value), so 95~\% of readings are within ±(30~ppm + 3~\%) & Whole ppm \\
Temperature & Gaussian, standard deviation 0.15~°C & Tenths of a degree \\
People counter & One too many in 5~\% of readings, one too few in 5~\% (never below zero) & A whole number, guards included \\
\bottomrule
\end{longtable}

Three faults can be set on one device (\code{deploy/fault.sh co2-0001 stuck}):
\emph{stuck} repeats the last good reading, \emph{dropout} reports nothing, and
\emph{drift} adds a rate per hour. A faulty reading is not marked in any way.

\subsection{6.5 What the model leaves out}

The model has the structure the control problem needs: people drive \COtwo{}
and heat, actuators answer with a lag, and the day has a shape. It is not a
building simulation. Each room is one well-mixed zone with no air or heat
exchanged with its neighbours or the corridor. The outdoor temperature and
\COtwo{} are constant, and there is no sun and no equipment load. Fan power is
taken as proportional to airflow, where a real fan's falls faster than that
when it slows down, so the proxy understates what low flows save. The thermal
constants are assumed, not fitted to any building. The policies in Section 11
are therefore compared against a model we wrote ourselves, and the result says
as much about that model's time constants as about the controllers.
